% !TEX root = ../main.tex
% Wave-11 W2/20 publication-ready inscription: N=1 (identity-twined) theorem.
% Author: Raeez Lorgat. No AI attribution anywhere. Atomic splice target.

\providecommand{\kBKM}{\kappa_{\mathrm{BKM}}}
\providecommand{\kch}{\kappa_{\mathrm{ch}}}
\providecommand{\kcat}{\kappa_{\mathrm{cat}}}
\providecommand{\kfib}{\kappa_{\mathrm{fiber}}}
\providecommand{\ClaimStatusTheorem}{\textbf{[Theorem]}}
\providecommand{\Hilb}{\mathrm{Hilb}}
\providecommand{\Mukai}{\Lambda_{\mathrm{Muk}}}
\providecommand{\gDeltaFive}{\mathfrak{g}_{\Delta_5}}

\section{The $N=1$ (identity-twined) theorem: $\gDeltaFive$, $\Phi_{10}$, and $K3 \times E$}
\label{sec:N1-theorem-publication}

The Igusa cusp form $\Phi_{10}$ is the square of Gritsenko's weight--$5$
Siegel form $\Delta_5$; $\Delta_5$ is the denominator of a rank--$3$
Borcherds--Kac--Moody superalgebra $\gDeltaFive$; the reduced Donaldson--Thomas
partition function of $K3 \times E$ is $-C/\Phi_{10}$. These three facts, each
proved in the primary literature, fit together into a single equality whose
left--hand side is a BKM root multiplicity, whose right--hand side is a BPS
count, and whose middle is the Fourier coefficient of a weak Jacobi form. The
theorem below states that equality.

\subsection{The identification}

Let $X = K3 \times E$ with $E$ an elliptic curve and $\Omega_X = \Omega_{K3} \wedge dz$
the holomorphic volume form. Write $\Mukai = H^\ast(K3, \bZ)$ for the Mukai
lattice (signature $(4,20)$); let $\Lambda^{2,1}_{\mathrm{II}} = \Lambda^{1,1} \oplus [2]$
be the rank--$3$ hyperbolic reflective sublattice that carries the real simple
roots $\{\delta_1, \delta_2, \delta_3\}$ of Gram matrix $\mathrm{diag}(2,2,2)$
with off--diagonal $-2$. Let $\phi_{0,1}$ denote the weak Jacobi form of weight
$0$ and index $1$ with expansion
\[
\phi_{0,1}(\tau, z) = (r^{-1} + 10 + r) + q\,(10 r^{-2} - 64 r^{-1} + 108 - 64 r + 10 r^2) + O(q^2),
\qquad r = e^{2\pi i z},\; q = e^{2\pi i \tau},
\]
and write $c(n, \ell)$ for the coefficient of $q^n r^\ell$. For a primitive
Mukai vector $\alpha = (n, \ell, m) \in \Lambda^{2,1}_{\mathrm{II}}$ with
$4nm - \ell^2 > 0$, denote by $\mathrm{mult}_{\gDeltaFive}(\alpha)$ the root
multiplicity of $\alpha$ in $\gDeltaFive$; for a primitive curve class
$\beta_h \in H_2(K3, \bZ)$ with $\beta_h^2 = 2h-2$ and a fibre degree
$\ell \in H_2(E, \bZ)$, denote by $n^0_{(\beta_h, \ell)}(X)$ the genus--zero
reduced BPS invariant of $X$ and by $N^{\mathrm{DT}}_{(\beta_h, \ell)}(X)$ the
corresponding reduced Donaldson--Thomas invariant.

\begin{theorem}[$N=1$: BKM--BPS--DT trifecta for $K3 \times E$]
\label{thm:N1-bkm-bps-dt}
\ClaimStatusTheorem\ Let $X = K3 \times E$ and let $\alpha = (n, \ell, m) \in
\Lambda^{2,1}_{\mathrm{II}}$ be a positive root of $\gDeltaFive$ with primitive
Mukai class corresponding under the rank--$3$ identification to
$(\beta_h, \ell) \in H_2(X, \bZ)$, $h = nm$. Then
\begin{equation}
\label{eq:N1-trifecta}
\mathrm{mult}_{\gDeltaFive}(\alpha) \;=\; c(nm, \ell) \;=\; n^0_{(\beta_h, \ell)}(X) \;=\; N^{\mathrm{DT}}_{(\beta_h, \ell)}(X),
\end{equation}
and the square of the BKM denominator is the reciprocal of the reduced DT
partition function up to the universal constant $C$ of Oberdieck--Pixton:
\begin{equation}
\label{eq:N1-Phi10-ZDT}
\Phi_{10}(\tau, z, \sigma) \;=\; \Delta_5(\tau, z, \sigma)^2 \;=\; -\,C \cdot Z^{\mathrm{red}}_{\mathrm{DT}}(X; q, y, \widetilde q)^{-1}.
\end{equation}
In particular the BKM weight invariant is
\[
\kBKM(\Delta_5) \;=\; \tfrac{1}{2}\, c_1(0) \;=\; \tfrac{1}{2} \cdot 10 \;=\; 5,
\]
realising the universal identity $\kBKM(\Phi_N) = c_N(0)/2$ at $N=1$.
\end{theorem}

\subsection{Proof}

The proof is a five--step chain. Each step is a published theorem; the content
of the present statement is that the four invariants in \eqref{eq:N1-trifecta}
coincide because each link in the chain pairs two of them.

\begin{proof}
\emph{Step 1: BKM side.} By Gritsenko--Nikulin \cite{GritsenkoNikulin1995, GritsenkoNikulin1998},
$\Delta_5$ is the denominator function of a generalised Kac--Moody superalgebra
$\gDeltaFive$ of rank $3$ with real simple root system $\{\delta_1, \delta_2, \delta_3\}$
supported on $\Lambda^{2,1}_{\mathrm{II}}$ and Weyl vector $\rho = f_2 - \tfrac12 f_3 + f_{-2}$.
Lemmas~2--3 of \emph{op.\,cit.} establish reflective anti--invariance under type
$\mathrm{II}$ reflections and invariance under type $\mathrm{I}$ reflections;
the Weyl--Kac--Borcherds denominator identity reads
\[
\Delta_5(Z) \;=\; e^{-2\pi i(\rho, Z)} \prod_{\alpha \in \Delta_+}\bigl(1 - e^{-2\pi i(\alpha, Z)}\bigr)^{\mathrm{mult}_{\gDeltaFive}(\alpha)}
\;=\; \sum_{w \in W^{(2)}}\mathrm{det}(w)\Bigl(e^{-2\pi i(w\rho, Z)} - \sum_a m(a)\, e^{-2\pi i (w(\rho + a), Z)}\Bigr).
\]
Reading off the exponent of $1 - e^{-2\pi i(\alpha, Z)}$ at
$\alpha = (n, \ell, m)$ gives $\mathrm{mult}_{\gDeltaFive}(\alpha) = c(nm, \ell)$.

\emph{Step 2: Borcherds lift and squaring.} Borcherds' singular theta lift
\cite[Thm~13.3]{Borcherds1998} sends $2\phi_{0,1}$ to a weight--$10$ Siegel
modular form; uniqueness of the Igusa cusp form yields
\[
\Phi_{10} \;=\; \mathrm{Bor}(2 \phi_{0,1}), \qquad \Delta_5^2 \;=\; \Phi_{10},
\]
so the squared denominator is the Igusa cusp form. The weight identity
$\kBKM(\Delta_5) = c_1(0)/2 = 5$ is the Borcherds weight formula applied to the
input $\phi_{0,1}$ (constant term $10$).

\emph{Step 3: DT side.} Oberdieck \cite{Oberdieck2018}, extending the
Oberdieck--Pandharipande programme, proves the Igusa cusp form conjecture: the
reduced equivariant Donaldson--Thomas partition function of $X = K3 \times E$
satisfies
\[
Z^{\mathrm{red}}_{\mathrm{DT}}(X;\, q, y, \widetilde q) \;=\; -\,\frac{C}{\Phi_{10}(\tau, z, \sigma)},
\]
where $q = e^{2\pi i \tau}$, $y = e^{2\pi i z}$, $\widetilde q = e^{2\pi i \sigma}$,
and $C$ is the Oberdieck--Pixton universal constant. This is the identity
\eqref{eq:N1-Phi10-ZDT}.

\emph{Step 4: BPS to Jacobi coefficient.} The Maulik--Pandharipande--Thomas
GW/PT correspondence on $K3$ \cite{MPT2010} together with the PT/DT correspondence
\cite{PT2014} and the Katz--Klemm--Vafa conjecture (proved for $K3$ by
Pandharipande--Thomas) identify the reduced genus--zero BPS count with the
Jacobi coefficient:
\[
n^0_{(\beta_h, \ell)}(X) \;=\; c(h, \ell) \;=\; c(nm, \ell).
\]
Integrality of the left--hand side is the content of Ionel--Parker
\cite{IonelParker2018}; equality to $c(h, \ell)$ follows from the Oberdieck--Pandharipande
multiplicative structure evaluated at $N = 1$ (the identity twist on $K3$).

\emph{Step 5: Wall--crossing to a single chamber.} Kontsevich--Soibelman
\cite{KS2008} and Joyce--Song \cite{JS2012} provide the product formula
$\Omega^{\mathrm{DT}} = \prod_\alpha (1 - q^\alpha)^{-a(\alpha)}$ for reduced
DT invariants on a CY$_3$ stratified by a stability condition. Evaluated at the
$\sigma_0$--chamber (the Oberdieck chamber for $K3 \times E$), the product
matches the Borcherds product expansion of $\Phi_{10}^{-1}$ term by term, so
\[
N^{\mathrm{DT}}_{(\beta_h, \ell)}(X) \;=\; n^0_{(\beta_h, \ell)}(X)
\]
for every primitive $(\beta_h, \ell)$.

Chaining Steps~1, 3, 4, 5 gives the four--way equality \eqref{eq:N1-trifecta};
Step~2 supplies \eqref{eq:N1-Phi10-ZDT} and the weight corollary. The
identification of the rank--$3$ BKM lattice with the primitive $(\beta_h, \ell)$
data on $K3 \times E$ is the $1 \leftrightarrow (\beta_h, \ell)$ bijection
$(n, \ell, m) \mapsto (\beta_{nm}, \ell)$.
\end{proof}

\subsection{Corollary on the $K3 \times E$ $\kappa_\bullet$--spectrum}

The $N=1$ identification contributes one of the four $\kappa_\bullet$ entries
in the $K3 \times E$ spectrum $\{2, 3, 5, 24\}$. The four entries come from
four distinct constructions: $\kcat(K3 \times E) = 0$ by the
K\"unneth--multiplicative rule (with the K3 fibre contributing $\chi(\cO_{K3}) = 2$
and $\chi(\cO_E) = 0$); $\kch(K3 \times E) = \sum_q (-1)^q h^{0,q}(K3 \times E) = 0$
by the Hodge--supertrace identification; $\kBKM(\Delta_5) = 5$ by
Theorem~\ref{thm:N1-bkm-bps-dt}; $\kfib(K3 \times E) = 24$ by the Euler number
of the K3 fibre under $\pi: K3 \times E \to E$. The four entries are not
permutations of one invariant computed four ways; they are four different
constructions, one of which is the BKM construction that Theorem~\ref{thm:N1-bkm-bps-dt}
pins down.

\subsection{Scope and dependencies}

The theorem is unconditional on the BKM side (Gritsenko--Nikulin, Borcherds),
unconditional on the DT side at $N = 1$ (Oberdieck's proof of the Igusa cusp
form conjecture), and unconditional on the BPS side for $K3$ fibre classes
(Pandharipande--Thomas KKV, Ionel--Parker integrality, MPT reduced GW/PT). The
only ingredient that requires a chamber choice is Step~5; the $\sigma_0$--chamber
is the one in which Oberdieck's partition function is defined, so the chamber
matches the chamber of Step~3 and no wall--crossing correction appears.

\begin{thebibliography}{99}
\bibitem{GritsenkoNikulin1995} V.~A.~Gritsenko, V.~V.~Nikulin,
\emph{Siegel automorphic form corrections of some Lorentzian Kac--Moody Lie algebras}, alg-geom/9504006, 1995.
\bibitem{GritsenkoNikulin1998} V.~A.~Gritsenko, V.~V.~Nikulin,
\emph{Automorphic forms and Lorentzian Kac--Moody algebras. Part II}, Internat. J. Math. 9 (1998), 201--275.
\bibitem{Borcherds1998} R.~E.~Borcherds,
\emph{Automorphic forms with singularities on Grassmannians}, Invent. Math. 132 (1998), 491--562.
\bibitem{Oberdieck2018} G.~Oberdieck,
\emph{Gromov--Witten invariants of the Hilbert schemes of points of a K3 surface}, Duke Math. J. 167 (2018), 3045--3125.
\bibitem{PT2014} R.~Pandharipande, R.~P.~Thomas,
\emph{The Katz--Klemm--Vafa conjecture for $K3$ surfaces}, Forum Math. Pi 4 (2016), e4.
\bibitem{MPT2010} D.~Maulik, R.~Pandharipande, R.~P.~Thomas,
\emph{Curves on $K3$ surfaces and modular forms}, J. Topol. 3 (2010), 937--996.
\bibitem{IonelParker2018} E.~Ionel, T.~H.~Parker,
\emph{The Gopakumar--Vafa formula for symplectic manifolds}, Ann. of Math. 187 (2018), 1--64.
\bibitem{KS2008} M.~Kontsevich, Y.~Soibelman,
\emph{Stability structures, motivic Donaldson--Thomas invariants and cluster transformations}, arXiv:0811.2435, 2008.
\bibitem{JS2012} D.~Joyce, Y.~Song,
\emph{A theory of generalized Donaldson--Thomas invariants}, Mem. Amer. Math. Soc. 217 (2012).
\end{thebibliography}
